% gen_event Handler Lifecycle Diagram
% Shows handler installation, event loop, swap, and supervised handler crash recovery.
% Compile with: xelatex gen_event_handler_lifecycle.tex
\documentclass[border=10pt]{standalone}

% ============================================================
% Packages
% ============================================================
\usepackage{fontspec}
\usepackage{xeCJK}
\setCJKmainfont[BoldFont=STHeiti, ItalicFont=STKaiti]{STSong}
\setCJKsansfont{STHeiti}
\setCJKmonofont{STFangsong}

\usepackage{xcolor}
\usepackage{tikz}
\usepackage{pifont}

\usetikzlibrary{arrows.meta, positioning, shapes.geometric, fit, calc, backgrounds, decorations.pathreplacing, decorations.markings}

% ============================================================
% Color Definitions
% ============================================================
\definecolor{erlred}{HTML}{A4070A}
\definecolor{erlblue}{HTML}{2B5EA7}
\definecolor{erlgreen}{HTML}{2E7D32}
\definecolor{erlyellow}{HTML}{F9A825}
\definecolor{erlpurple}{HTML}{6A1B9A}
\definecolor{erlmodule}{HTML}{D84315}
\definecolor{erlinit}{HTML}{00897B}
\definecolor{erlpid}{HTML}{E65100}
\definecolor{erlteal}{HTML}{00695C}
\definecolor{erlhandler}{HTML}{1565C0}
\definecolor{erlevent}{HTML}{E65100}
\definecolor{erlmanager}{HTML}{283593}
\definecolor{erlswap}{HTML}{6A1B9A}
\definecolor{erlcrash}{HTML}{C62828}
\definecolor{darkgray}{HTML}{333333}

% ============================================================
% Document
% ============================================================
\begin{document}

\begin{tikzpicture}[scale=0.72, every node/.style={scale=0.72},
    phase/.style={draw=#1!70, fill=#1!10, rounded corners=4pt,
                  minimum width=2.8cm, minimum height=1.4cm,
                  font=\small\bfseries, align=center},
    arrow/.style={-{Stealth[length=5pt]}, thick, #1, line width=1.2pt},
]

    % ==================================================================
    % PART 1: Normal Handler Lifecycle (top section)
    % ==================================================================
    \node[font=\large\bfseries\color{erlblue}, anchor=west] at (-8, 7.5)
        {Handler 生命周期};

    % --- Phase nodes ---
    \node[phase=erlgreen] (add) at (-6, 5.5) {add\_handler\\{\tiny\texttt{(Mgr, Mod, Args)}}};
    \node[phase=erlinit] (init) at (-2, 5.5) {Mod:init/1\\{\tiny\texttt{\{ok, State\}}}};
    \node[phase=erlblue] (loop) at (3, 5.5) {事件循环\\{\tiny handle\_event/2}\\{\tiny handle\_call/2}\\{\tiny handle\_info/2}};
    \node[phase=erlred] (term) at (8, 5.5) {Mod:terminate/2\\{\tiny 清理资源}};

    % --- Arrows ---
    \draw[arrow=erlgreen!80!black] (add) -- node[above, font=\tiny\ttfamily] {调用回调} (init);
    \draw[arrow=erlinit!80!black] (init) -- node[above, font=\tiny] {进入循环} (loop);
    \draw[arrow=erlred!80!black] (loop) -- node[above, font=\tiny] {delete / stop / 崩溃} (term);

    % Self-loop on event loop
    \draw[arrow=erlblue!80!black] (loop.north) .. controls +(0,1.2) and +(0,1.2) ..
        node[above, font=\tiny, align=center] {Event / Call / Info\\{\tiny 更新 State}} (loop.north east);

    % --- Return values annotation ---
    \node[draw=erlblue!40, fill=erlblue!5, rounded corners=3pt,
          font=\tiny\ttfamily, text width=6cm, align=left, anchor=north west]
        at (-6, 3.5) {
        handle\_event(Event, State) ->\\
        \quad \{ok, NewState\}\\
        \quad | \{ok, NewState, hibernate\}\\
        \quad | \{swap\_handler, Args1, NewState, NewHandler, Args2\}\\
        \quad | remove\_handler\\[3pt]
        handle\_call(Request, State) ->\\
        \quad \{ok, Reply, NewState\}\\
        \quad | \{ok, Reply, NewState, hibernate\}\\
        \quad | \{swap\_handler, Reply, Args1, NewState, NewHandler, Args2\}\\
        \quad | \{remove\_handler, Reply\}
    };

    % ==================================================================
    % PART 2: swap_handler Flow (middle section)
    % ==================================================================
    \node[font=\large\bfseries\color{erlswap}, anchor=west] at (-8, 1) {swap\_handler 热替换};

    % Old handler
    \node[draw=erlswap!70, fill=erlswap!10, rounded corners=5pt,
          minimum width=3.5cm, minimum height=2cm] (old) at (-4.5, -1) {};
    \node[font=\small\bfseries\color{erlswap}, anchor=north] at (old.north) {旧 Handler};
    \node[font=\tiny\ttfamily, align=center] at (old.center) {
        Module: old\_handler\\
        State: OldState
    };

    % Terminate box
    \node[draw=erlred!60, fill=erlred!8, rounded corners=3pt,
          minimum width=2.5cm, minimum height=1.2cm,
          font=\scriptsize\ttfamily, align=center] (oldterm) at (-0.5, -1) {
        terminate(\\SwapArgs,\\OldState)
    };

    % Transfer data
    \node[draw=erlyellow!80, fill=erlyellow!15, rounded corners=3pt,
          minimum width=2.2cm, minimum height=0.9cm,
          font=\scriptsize\bfseries, align=center] (transfer) at (3, -1) {
        TransferData\\
        {\tiny (返回值)}
    };

    % New handler init
    \node[draw=erlinit!60, fill=erlinit!8, rounded corners=3pt,
          minimum width=2.8cm, minimum height=1.2cm,
          font=\scriptsize\ttfamily, align=center] (newinit) at (6.5, -1) {
        new\_handler:\\
        init(\{Args,\\TransferData\})
    };

    % New handler
    \node[draw=erlgreen!70, fill=erlgreen!10, rounded corners=5pt,
          minimum width=3.5cm, minimum height=2cm] (new) at (10, -1) {};
    \node[font=\small\bfseries\color{erlgreen!80!black}, anchor=north] at (new.north) {新 Handler};
    \node[font=\tiny\ttfamily, align=center] at (new.center) {
        Module: new\_handler\\
        State: NewState
    };

    % Arrows for swap
    \draw[arrow=erlswap] (old) -- (oldterm);
    \draw[arrow=erlyellow!80!black] (oldterm) -- node[above, font=\tiny] {返回} (transfer);
    \draw[arrow=erlinit!80!black] (transfer) -- node[above, font=\tiny] {传递} (newinit);
    \draw[arrow=erlgreen!80!black] (newinit) -- (new);

    % Brace annotation
    \draw[decorate, decoration={brace, amplitude=6pt, mirror}, thick, erlswap!60]
        ([yshift=-8pt]old.south west) -- ([yshift=-8pt]new.south east)
        node[midway, below=8pt, font=\tiny\color{erlswap!80!black}, align=center] {
            整个过程在同一个 Manager 进程中完成\\
            Handler 列表中旧条目被新条目\textbf{原地替换}
        };

    % ==================================================================
    % PART 3: add_sup_handler & Crash Recovery (bottom section)
    % ==================================================================
    \node[font=\large\bfseries\color{erlcrash}, anchor=west] at (-8, -4)
        {add\_sup\_handler 监督模式};

    % Supervisor process
    \node[draw=erlgreen!70, fill=erlgreen!8, rounded corners=5pt,
          minimum width=3.5cm, minimum height=2.2cm] (sup) at (-5, -6.5) {};
    \node[font=\small\bfseries\color{erlgreen!80!black}, anchor=north] at (sup.north)
        {监督进程 (Caller)};
    \node[font=\tiny\ttfamily, align=center] at ([yshift=-2pt]sup.center) {
        add\_sup\_handler(Mgr,\\
        \quad Handler, Args)\\[2pt]
        {\tiny 自动 monitor Handler}
    };

    % Event Manager
    \node[draw=erlmanager!70, fill=erlmanager!8, rounded corners=5pt,
          minimum width=3.5cm, minimum height=2.2cm] (mgr) at (1, -6.5) {};
    \node[font=\small\bfseries\color{erlmanager}, anchor=north] at (mgr.north)
        {Event Manager};
    \node[font=\tiny\ttfamily, align=center] at ([yshift=-2pt]mgr.center) {
        Handler 正常运行中...\\[4pt]
        {\tiny\textcolor{erlcrash}{\textbf{Handler 崩溃!}}}\\
        {\tiny 从列表中移除}
    };

    % Crash notification
    \node[draw=erlcrash!70, fill=erlcrash!8, rounded corners=4pt,
          minimum width=4cm, minimum height=1.6cm] (crash) at (7, -6.5) {};
    \node[font=\small\bfseries\color{erlcrash}, anchor=north] at (crash.north)
        {崩溃通知};
    \node[font=\tiny\ttfamily, align=center] at ([yshift=-2pt]crash.center) {
        \{gen\_event\_EXIT,\\
        \quad Handler, Reason\}\\[2pt]
        {\tiny 发送给监督进程}
    };

    % Arrows
    \draw[arrow=erlgreen!80!black] (sup.east) --
        node[above, font=\tiny, align=center] {\ding{172} 安装\\{\tiny + link}} (mgr.west);
    \draw[arrow=erlcrash] (mgr.east) --
        node[above, font=\tiny, align=center] {\ding{173} 崩溃} (crash.west);
    \draw[arrow=erlcrash, bend right=30] (crash.south) to
        node[below, font=\tiny, align=center] {\ding{174} 通知监督进程\\{\tiny 可重新安装 Handler}}
        (sup.south);

    % --- Key difference box ---
    \node[draw=erlcrash!50, fill=erlcrash!5, rounded corners=3pt,
          font=\tiny, text width=14cm, align=left, anchor=north west]
        at (-8, -8.8) {
        \textbf{add\_handler vs add\_sup\_handler 的区别：}\\[2pt]
        \texttt{add\_handler}：Handler 崩溃后静默移除，无人知晓\\
        \texttt{add\_sup\_handler}：Handler 崩溃后向调用者发送 \texttt{\{gen\_event\_EXIT, Handler, Reason\}} 消息，
        调用者可据此重新安装 Handler 或采取其他恢复措施
    };

\end{tikzpicture}

\end{document}
